\documentclass[CJK,13pt]{beamer}
\usepackage{CJKutf8}
\usepackage{beamerthemesplit}
\usetheme{Malmoe}
\useoutertheme[footline=authortitle]{miniframes}
\usepackage{amsmath}
\usepackage{amssymb}
\usepackage{graphicx}
\usepackage{eufrak}
\usepackage{color}
\usepackage{slashed}
\usepackage{simplewick}
\usepackage{tikz}
\usepackage{tcolorbox}
\usepackage{ulem}
\graphicspath{{../figures/}}
%%figures
\def\lfig#1#2{\includegraphics[width=#1 in]{#2}}
\def\tfig#1#2{\includegraphics[height=#1 in]{#2}}
\def\addfig#1#2{\begin{center}\includegraphics[width=#1 in]{#2}\end{center}}
\def\question{\includegraphics[width=0.3in]{why.jpg}\,}
\def\answer{\includegraphics[width=0.3in]{baozou_haha.png}\,}
\def\wulian{\includegraphics[width=0.18in]{emoji_wulian.jpg}}
\def\bigwulian{\includegraphics[width=0.35in]{emoji_wulian.jpg}}
\def\bye{\includegraphics[width=0.18in]{emoji_bye.jpg}}
\def\bigbye{\includegraphics[width=0.35in]{emoji_bye.jpg}}
\def\huaixiao{\includegraphics[width=0.18in]{emoji_huaixiao.jpg}}
\def\bighuaixiao{\includegraphics[width=0.35in]{emoji_huaixiao.jpg}}
\def\jianxiao{\includegraphics[width=0.18in]{emoji_jianxiao.jpg}}
\def\bigjianxiao{\includegraphics[width=0.35in]{emoji_jianxiao.jpg}}
\def\haoqi{\includegraphics[width=0.18in]{emoji_haoqi.jpg}}
%% colors
\def\blacktext#1{{\color{black}#1}}
\def\bluetext#1{{\color{blue}#1}}
\def\redtext#1{{\color{red}#1}}
\def\darkbluetext#1{{\color[rgb]{0,0.2,0.6}#1}}
\def\skybluetext#1{{\color[rgb]{0.2,0.7,1.}#1}}
\def\cyantext#1{{\color[rgb]{0.,0.5,0.5}#1}}
\def\greentext#1{{\color[rgb]{0,0.7,0.1}#1}}
\def\darkgray{\color[rgb]{0.2,0.2,0.2}}
\def\lightgray{\color[rgb]{0.6,0.6,0.6}}
\def\gray{\color[rgb]{0.4,0.4,0.4}}
\def\blue{\color{blue}}
\def\red{\color{red}}
\def\orange{\color[rgb]{1.,0.8,0.}}
\def\green{\color{green}}
\def\darkgreen{\color[rgb]{0,0.4,0.1}}
\def\darkblue{\color[rgb]{0,0.2,0.6}}
\def\skyblue{\color[rgb]{0.2,0.7,1.}}
%%control
\def\diag{\mathrm{diag}\,}
\def\heaviside{\,\mathrm{h}\,}
\def\bral{\left(\begin{array}{l}}
\def\brar{\end{array}\right)}
\def\brall{\left(\begin{array}{ll}}
\def\brarr{\end{array}\right)}
\def\bralll{\left(\begin{array}{lll}}
\def\brarrr{\end{array}\right)}
\def\branchl{\left\{\begin{array}{l}}
\def\branchr{\end{array}\right.}
\def\branchll{\left\{\begin{array}{ll}}
\def\branchrr{\end{array}\right.}
\def\branchlll{\left\{\begin{array}{lll}}
\def\branchrrr{\end{array}\right.}
\def\sfgamma#1{\,\Gamma\left( #1 \right)\,}
\def\be{\begin{equation}}
\def\ee{\nonumber\end{equation}}
\def\bea{\begin{eqnarray}}
\def\eea{\nonumber\end{eqnarray}}
\def\bch{\begin{CJK}{UTF8}{gbsn}}
\def\ech{\end{CJK}}
\def\bitem{\begin{itemize}}
\def\eitem{\end{itemize}}
\def\bcenter{\begin{center}}
\def\ecenter{\end{center}}
\def\bex{\begin{minipage}{0.2\textwidth}\includegraphics[width=0.6in]{jugelizi.png}\end{minipage}\begin{minipage}{0.76\textwidth}}
\def\eex{\end{minipage}}
\def\chtitle#1{\frametitle{\bch#1\ech}}
\def\bmat#1{\left(\begin{array}{#1}}
\def\emat{\end{array}\right)}
\def\bcase#1{\left\{\begin{array}{#1}}
\def\ecase{\end{array}\right.}
\def\bmini#1{\begin{minipage}{#1\textwidth}}
\def\emini{\end{minipage}}
\def\tbox#1{\begin{tcolorbox}#1\end{tcolorbox}}
\def\pfrac#1#2#3{\left(\frac{\partial #1}{\partial #2}\right)_{#3}}
\def\res#1#2{\,\mathrm{res}\,#1\left(#2\right)\,}
\def\newt#1#2{\left(\begin{array}{c}#1\\ #2\end{array}\right)}
\def\reof#1{\,\mathrm{Re}{\left(#1\right)\,}}
\def\imof#1{\,\mathrm{Im}{\left(#1\right)\,}}
\def\Arg#1{\,\mathrm{Arg}\,#1\,}
%%symbols
\def\intfull{\int_{-\infty}^\infty}
\def\inthalf{\int_0^\infty}
\def\bropt{\,(\ \ \ )}
\def\sech{\mathrm{sech}\,}
\def\csch{\mathrm{csch}\,}
\def\asinh{\mathrm{asinh}\,}
\def\acosh{\mathrm{acosh}\,}
\def\sone{$\star$}
\def\stwo{$\star\star$}
\def\sthree{$\star\star\star$}
\def\sfour{$\star\star\star\star$}
\def\sfive{$\star\star\star\star\star$}
\def\rint{{\int_\leftrightarrow}}
\def\roint{{\oint_\leftrightarrow}}
\def\stdHf{{\textit{\r H}_f}}
\def\deltaH{{\Delta \textit{\r H}}}
\def\ii{{\dot{\imath}}}
\def\skipline{{\vskip0.1in}}
\def\skiplines{{\vskip0.2in}}
\def\lagr{{\mathcal{L}}}
\def\hamil{{\mathcal{H}}}
\def\vecT{{\mathbf{T}}}
\def\vecN{{\mathbf{N}}}
\def\vecB{{\mathbf{B}}}
\def\vecv{{\mathbf{v}}}
\def\vecr{{\mathbf{r}}}
\def\vecf{{\mathbf{f}}}
\def\vecg{{\mathbf{g}}}
\def\vecupsilon{{\mathbf{\upsilon}}}
\def\vecu{{\mathbf{u}}}
\def\vecj{{\mathbf{j}}}
\def\vecx{{\mathbf{x}}}
\def\vecy{{\mathbf{y}}}
\def\vecz{{\mathbf{z}}}
\def\veck{{\mathbf{k}}}
\def\vecp{{\mathbf{p}}}
\def\vecn{{\mathbf{n}}}
\def\vecA{{\mathbf{A}}}
\def\vecP{{\mathbf{P}}}
\def\vecsigma{{\mathbf{\sigma}}}
\def\hatJn{{\hat{J_\vecn}}}
\def\hatJx{{\hat{J_x}}}
\def\hatJy{{\hat{J_y}}}
\def\hatJz{{\hat{J_z}}}
\def\hatj#1{\hat{J_{#1}}}
\def\hatphi{{\hat{\phi}}}
\def\hatq{{\hat{q}}}
\def\hatpi{{\hat{\pi}}}
\def\vel{\upsilon}
\def\Dint{{\mathcal{D}}}
\def\adag{{\hat{a}^\dagger}}
\def\bdag{{\hat{b}^\dagger}}
\def\cdag{{\hat{c}^\dagger}}
\def\ddag{{\hat{d}^\dagger}}
\def\hata{{\hat{a}}}
\def\hatb{{\hat{b}}}
\def\hatc{{\hat{c}}}
\def\hatd{{\hat{d}}}
\def\hatD{{\,\hat{D}}}
\def\hatN{{\hat{N}}}
\def\hatH{{\hat{H}}}
\def\hatp{{\hat{p}}}
\def\Fup{{F^{\mu\nu}}}
\def\Fdown{{F_{\mu\nu}}}
\def\newl{\nonumber \\}
\def\vece{\mathrm{e}}
\def\calM{{\mathcal{M}}}
\def\calT{{\mathcal{T}}}
\def\calR{{\mathcal{R}}}
\def\barpsi{\bar{\psi}}
\def\baru{\bar{u}}
\def\barv{\bar{\upsilon}}
\def\qeq{\stackrel{?}{=}}
\def\ftf{\stackrel{\mathcal{FT}}{\Longrightarrow}}
\def\ftb{\stackrel{\mathcal{FT}}{\Longleftarrow}}
\def\ftfb{\stackrel{\mathcal{FT}}{\Longleftrightarrow}}
\def\ltf{\stackrel{\mathcal{LT}}{\Longrightarrow}}
\def\ltb{\stackrel{\mathcal{LT}}{\Longleftarrow}}
\def\ltfb{\stackrel{\mathcal{LT}}{\Longleftrightarrow}}
\def\torder#1{\mathcal{T}\left(#1\right)}
\def\rorder#1{\mathcal{R}\left(#1\right)}
\def\contr#1#2{\contraction{}{#1}{}{#2}#1#2}
\def\trof#1{\mathrm{Tr}\left(#1\right)}
\def\trace{\mathrm{Tr}}
\def\comm#1{\ \ \ \left(\mathrm{used}\ #1\right)}
\def\tcomm#1{\ \ \ (\text{#1})}
\def\slp{\slashed{p}}
\def\slk{\slashed{k}}
\def\calp{{\mathfrak{p}}}
\def\veccalp{\mathbf{\mathfrak{p}}}
\def\Tthree{T_{\tiny \textcircled{3}}}
\def\pthree{p_{\tiny \textcircled{3}}}
\def\dbar{{\,\mathchar'26\mkern-12mu d}}
\def\erf{\mathrm{erf}}
\def\const{\mathrm{const.}}
\def\pheat{\pfrac p{\ln T}V}
\def\vheat{\pfrac V{\ln T}p}
%%units
\def\fdeg{{^\circ \mathrm{F}}}
\def\cdeg{^\circ \mathrm{C}}
\def\atm{\,\mathrm{atm}}
\def\angstrom{\,\text{\AA}}
\def\SIL{\,\mathrm{L}}
\def\SIkm{\,\mathrm{km}}
\def\SIyr{\,\mathrm{yr}}
\def\SIGyr{\,\mathrm{Gyr}}
\def\SIV{\,\mathrm{V}}
\def\SImV{\,\mathrm{mV}}
\def\SIeV{\,\mathrm{eV}}
\def\SIkeV{\,\mathrm{keV}}
\def\SIMeV{\,\mathrm{MeV}}
\def\SIGeV{\,\mathrm{GeV}}
\def\SIcal{\,\mathrm{cal}}
\def\SIkcal{\,\mathrm{kcal}}
\def\SImol{\,\mathrm{mol}}
\def\SIN{\,\mathrm{N}}
\def\SIHz{\,\mathrm{Hz}}
\def\SIm{\,\mathrm{m}}
\def\SIcm{\,\mathrm{cm}}
\def\SIfm{\,\mathrm{fm}}
\def\SImm{\,\mathrm{mm}}
\def\SInm{\,\mathrm{nm}}
\def\SImum{\,\mathrm{\mu m}}
\def\SIJ{\,\mathrm{J}}
\def\SIW{\,\mathrm{W}}
\def\SIkJ{\,\mathrm{kJ}}
\def\SIs{\,\mathrm{s}}
\def\SIkg{\,\mathrm{kg}}
\def\SIg{\,\mathrm{g}}
\def\SIK{\,\mathrm{K}}
\def\SImmHg{\,\mathrm{mmHg}}
\def\SIPa{\,\mathrm{Pa}}
%page
\def\secpage#1#2{\begin{frame}\bcenter{\bf \Huge #1} \ecenter \skipline \tbox{\bcenter #2 \ecenter}\end{frame}}
\def\append#1#2{\secpage{附录#1}{#2}}
\def\thinka#1{\begin{frame}\frametitle{思考题}\bcenter\lfig{0.4}{think0.jpg}  \skipline #1 \ecenter\end{frame}}
\def\thinkb#1{\begin{frame}\frametitle{思考题}\bcenter\lfig{0.5}{think1.jpg}  \skipline #1 \ecenter\end{frame}}
\def\thinkc#1{\begin{frame}\frametitle{思考题}\bcenter\lfig{0.5}{think2.jpg}  \skipline #1 \ecenter\end{frame}}
\def\thinkd#1{\begin{frame}\frametitle{思考题}\bcenter\lfig{0.5}{think3.jpg}  \skipline #1 \ecenter\end{frame}}
\def\thinke#1{\begin{frame}\frametitle{思考题}\bcenter\lfig{0.5}{think4.jpg}  \skipline #1 \ecenter\end{frame}}
\def\thinkf#1{\begin{frame}\frametitle{思考题}\bcenter\lfig{0.5}{think5.jpg}  \skipline #1 \ecenter\end{frame}}
\def\schw{ds^2 = \left(1-\frac{2GM}{r}\right)dt^2 - \left(1-\frac{2GM}{r}\right)^{-1} dr^2 - r^2\left(d\theta^2 + \sin^2\theta d\phi^2\right)}

\input{titlepage.tex}
  \date{}
  \begin{document}
  \bch
  \tpage{21}{Physics of Gravitational Waves}

  \begin{frame}
    \frametitle{专业版引力波三连锤}
    上一讲介绍了科普版引力波三连锤。这一讲我们介绍专业版：
    \bitem
  \item{引力（波）只有2个物理自由度，神马意思？}
  \item{引力波是横向无迹的，解释下？}        
  \item{为什么引力波自旋是2？}
    \eitem

    \addfig{3}{shanchu.jpg}
  \end{frame}
  
  \secpage{引力（波）的物理自由度？}{只有2个}
  
  \begin{frame}
    \frametitle{引力波是物理的吗？}
    我们约定了谐和坐标条件
    $$ \square x^\mu = 0$$
    然后导出了引力波要满足的爱因斯坦方程
    $$ \square h_{\mu\nu} = 0$$
    是不是忍不住要怀疑 $h_{\mu\nu}$ 是人为操作出来的？
  \end{frame}



  \begin{frame}
    注意到谐和坐标系不是唯一的，对 $\square \epsilon^\mu = 0$ （学过电动力学的你丝毫不会怀疑这样的 $\epsilon^\mu$ 有无穷多种解）且 $\epsilon^\mu$ 不超过 $h_{\mu\nu}$ 的量级，那么至少在一阶近似下，坐标变换 $x^\mu \rightarrow x^\mu+\epsilon^\mu$ 仍然给出谐和坐标系。

    \skipline
    
    看来操作的空间很大——

    \skipline
    
    愈发怀疑……
  \end{frame}

  \begin{frame}
    为了简化讨论，我们来考虑一个沿 $z$ 轴传播的平面波解。
    \bitem
  \item{$h_{33}$ 描述的是沿传播方向的度量（即物理长度和$z$坐标之间的关系）的振动，这家伙可以通过操作 $z$ 坐标的定义来消除。}
  \item{虽然这有些不太好想象，但是数学上的对称性容易让你相信操作时间坐标可以干掉 $h_{00}$。}
  \item{虽然引力波并不在 $x$ 和 $y$ 方向上移动，但是操作 $x$ 坐标和 $y$ 坐标还是可以干掉 $h_{13}$ 和 $h_{23}$（这类似于动态地调节过$z$轴上各点的 $x$, $y$ 方向和 $z$轴的夹角来消除度规交叉项）。}
    \eitem
  \end{frame}

  \begin{frame}
    注意傅立叶波矢  $k^\mu= (\omega, 0, 0, \omega)$, $k_\mu = (\omega, 0, 0, -\omega)$，谐和坐标条件
    $$\eta^{\mu\nu}h_{\mu\alpha,\nu} = \frac{1}{2}h_{,\alpha}$$
    要求
    $$ k^\mu h_{\mu\alpha} = \frac{1}{2}k_\alpha h$$
    令 $\alpha=0,1,2,3$，分别得到
    $$  h_{30} = \frac{1}{2}h, h_{01}=0, h_{02}=0, h_{03}=-\frac{1}{2}h $$
    所以，$h_{01},h_{02},h_{03}$也自动消失了，且 $h_{11}+h_{22}=0$。

    \skipline
    
     这时候 $h_{\mu\nu}$ 的迹 $h=0$，且满足 $k^\mu h_{\mu\alpha} = 0$，这叫做{\blue 横向无迹(Transverse-Traceless，简称TT)规范}。
  \end{frame}


  \begin{frame}

    最后的结论是： {\blue 对沿 $z$ 轴方向传播的平面波，只有 $h_{11}-h_{22}$ 和 $h_{12}$ 是物理的两个自由度。}其余所有分量以及 $h_{11}+h_{22}$ 都是零。

    \skiplines

  \end{frame}

  \begin{frame}
    
    度规有10个自由度，四个坐标选择的随意性干掉了4个非物理自由度。那么度规应该还有6个物理自由度才对。

    \skiplines
    
    有时候我们会说，引力（没有波字！）只有两个物理自由度。这是什么意思呢？

    \skiplines

    因为引力波可以脱离物质自己传播，所以对应引力的内禀自由度。而度规的其他4个物理自由度由类似于牛顿引力的泊松方程描述（幼儿版GR就不进行这些繁琐的数学推导了），拿掉物质的源，响应就会消失。所以可以认为度规的其他物理自由度其实是物质的自由度，而不是引力本身的。
  \end{frame}

  \secpage{引力子}{是自旋为2的粒子}

  \begin{frame}
    我们经常在科普读物中看到介绍说引力子的自旋为$2$，这是啥意思呢？
  \end{frame}

  \begin{frame}
    我们都知道平面上自旋为 $1$ 的，也就是矢量场 $\vecv$ 在坐标系旋转 $\theta$  角时，满足下列变换规则:
    \be
      \tilde{v}_x = v_x\cos{\theta} + v_y\sin{\theta},\ \tilde{v}_y = v_y\cos{\theta} - v_x\sin{\theta}.
    \ee
    自旋为2的场的变换规则很简单，就是把 $\theta$ 变为 $2\theta$。
    
    \addfig{1.}{chulecai.jpg}

    但是——这好像有些超出本喵的想象力
  \end{frame}
  
  
  \begin{frame}
    为了获得一些直观理解，我们先来看一些由$x$-$y$平面上的标量场构造出来的自旋为 $2$ 的场。设 $E$ 是一个标量场，定义
    \be
      q_E \equiv \frac{1}{2}\left(\partial_x^2 - \partial_y^2\right)E,\ u_E = \partial_x\partial_yE.
    \ee
  \end{frame}

  \begin{frame}
    当$x$-$y$坐标系绕原点旋转$\theta$至 $\tilde{x}$-$\tilde{y}$坐标系时，
    \be
      \frac{\partial x}{\partial\tilde{x}} = \cos\theta,\ \frac{\partial x}{\partial \tilde{y}} = -\sin\theta, \ \frac{\partial y}{\partial\tilde{x}} = \sin\theta,\ \frac{\partial y}{\partial \tilde{y}} = \cos\theta, \ 
      \ee
    于是有
    \be
      \frac{\partial E}{\partial\tilde{x}} = \frac{\partial E}{\partial x} \cos\theta + \frac{\partial E}{\partial y}\sin\theta,\ \frac{\partial E}{\partial \tilde{y}} = \frac{\partial E}{\partial y}\cos\theta-  \frac{\partial E}{\partial x} \sin\theta.
    \ee      
  \end{frame}

  \begin{frame}
    以及
    \begin{eqnarray}
      \frac{\partial^2 E}{\partial\tilde{x}^2} &=& \frac{\partial^2 E}{\partial x^2} \cos^2\theta + \frac{\partial^2 E}{\partial y^2}\sin^2\theta + \frac{\partial^2 E}{\partial x \partial y}\sin 2\theta, \\
      \frac{\partial^2 E}{\partial\tilde{x}\partial\tilde{y}} &=& \frac{\partial^2E}{\partial x \partial y} \cos 2\theta  - \frac{1}{2}\left(\frac{\partial^2 E}{\partial x^2}-\frac{\partial^2 E}{\partial y^2}\right)\sin 2\theta, \\  
      \frac{\partial^2 E}{\partial\tilde{y}^2} &=& \frac{\partial^2E}{\partial y^2} \cos^2\theta + \frac{\partial^2E}{\partial x^2}\sin^2\theta - \frac{\partial^2E}{\partial x \partial y}\sin{2\theta},   
    \end{eqnarray}
    或者写成更紧凑的形式
    \be
      \tilde{q}_E = q_E\cos{2\theta} + u_E\sin{2\theta},\ \tilde{u}_E = u_E\cos{2\theta} - q_E\sin{2\theta}.
    \ee
    这说明我们从标量场 $E$ 构造的 $(q_E, u_E)$ 是自旋为 $2$ 的场。
  \end{frame}

  \begin{frame}
    对于沿 $z$ 轴方向传播($x,y$坐标均为零）的引力波，假设我们已经操作掉了其他自由度，只剩下 $h_{11}-h_{22}$ 和 $h_{12}$。

      对局域的“标量”函数
      $$ E(x,y) = \frac{1}{2}h_{11}x^2 +  \frac{1}{2}h_{22}y^2 + h_{12}xy $$
      （它只是局部的 $\frac{1}{2}ds^2$ 而已)
        计算 $q_E, u_E$， 即可知道 $\left(\frac{h_{11}-h_{22}}{2}, h_{12}\right)$ 是自旋为2的场。
        在文献中，一般不写 $q_E, u_E$，而是用符号
        $$h_+ \equiv \frac{h_{11}-h_{22}}{2}, h_\times \equiv h_{12}$$
        来表示引力波的两种模式。
  \end{frame}

  \begin{frame}
    \addfig{3}{hpluscross.png}
  \end{frame}

  
  \begin{frame}
    要注意的是，$ds^2$ 不能是全局的标量，因此引力波(微扰意义下的二阶张量）不能用一个全局的标量场来描述。
  \end{frame}

  
    \ech
\end{document}




  
